\chapter{Limitations}\label{ch:limitations}

\section{No Quantitative Evaluation}

The most significant limitation of this project is the complete absence of systematic evaluation. We have not constructed a labelled benchmark, measured retrieval precision or recall, or produced any quantitative scores for answer faithfulness or relevance.

No evaluation dataset for SUSTech-specific information retrieval exists. Constructing one would require: (a) formulating a representative set of questions covering the full range of query types (factual lookups, procedural queries, multi-document synthesis, and out-of-scope rejections); (b) manually verifying ground-truth answers against the authoritative source pages; and (c) ensuring the crawled corpus actually contains those source pages. This is non-trivial work, and skipping it means the design choices in this project---paragraph-first chunking, hybrid retrieval, the self-RAG loop---are motivated by prior literature and engineering reasoning, but their actual effect on answer quality over this specific corpus is unknown.

Without a benchmark, it is impossible to run controlled ablations. We cannot quantify the contribution of BM25 over dense-only retrieval, the improvement from cross-encoder reranking, or the difference in answer correctness between simple RAG and self-RAG modes. The system is a working prototype, but not a validated one.

Building a proper evaluation set---even a modest 100--200 labelled (question, answer, source passage) triples---would be the single most impactful next step. It would allow ablations across retrieval parameters, chunk sizes, and self-RAG configurations, and would transform the system from a well-engineered prototype into one with measurable, bounded properties.

\section{Data Coverage and Freshness}

The knowledge base was built from a static crawl limited to approximately 200 pages. This is a small fraction of the full SUSTech web presence. Several categories of pages are systematically absent:

\paragraph{JavaScript-rendered content.} The crawler operates at the HTML level and cannot execute JavaScript. Pages that load their content dynamically (staff directories, event listings, course catalogues) are either empty or absent from the index after extraction. This includes some of the most frequently queried types of campus information.

\paragraph{Data freshness.} Campus information changes on short timescales: staff contacts, course offerings, event notices, and administrative procedures are updated frequently. The current system has no incremental update mechanism. A scheduled re-crawl with page-level change detection would be required to keep the knowledge base current.

\paragraph{Crawl depth and seed coverage.} The BFS crawl starts from a fixed set of seed URLs. Subsites that are not reachable from those seeds---or that are reachable only through JavaScript navigation---are never discovered. Important information sources such as departmental faculty pages, the graduate admissions portal, and the library catalogue are likely underrepresented or absent.

Together, these gaps mean that many legitimate campus queries will result in a ``not found in knowledge base'' response, not because the system failed to retrieve correctly, but because the relevant page was never indexed. Expanding the crawl, adding targeted seed URLs, and integrating a headless browser for JavaScript-rendered pages would substantially improve coverage.
